% !TEX program = xelatex
\documentclass[11pt]{article}

\usepackage[a4paper,margin=1in]{geometry}
\usepackage{amsmath,amssymb}
\usepackage{booktabs}
\usepackage{enumitem}
\usepackage{hyperref}
\usepackage{xcolor}
\usepackage{kotex}
\usepackage{listings}

\lstdefinestyle{codestyle}{
    basicstyle=\ttfamily\small,
    breaklines=true,
    frame=single,
    columns=fullflexible,
    backgroundcolor=\color{gray!5}
}
\lstset{style=codestyle}

\title{FDM vs PINN Full-Field Comparison Notebook Guide\\
\large \texttt{fdm\_vs\_pinn\_full\_field\_comparison.ipynb} 셀별 코드 설명}
\author{안형준}
\date{}

\begin{document}
\maketitle

\section{문서 목적}

이 문서는 \texttt{fdm\_vs\_pinn\_full\_field\_comparison.ipynb}의 각 셀이 무슨 역할을 하는지 설명한다.
노트북의 목적은 finite difference method(FDM)로 만든 기준 해와 DeepXDE PINN 예측 해를 같은 \(X-\tau\) 전체 격자에서 비교하는 것이다.

\[
\theta_{\mathrm{FDM}}(X,\tau)
\quad \text{vs.} \quad
\theta_{\mathrm{PINN}}(X,\tau)
\]

비교 대상은 sensor 위치의 값만이 아니라 전체 temperature field이다.
따라서 이 노트북은 단순히 학습 결과의 \(\beta\), \(S\), \(h\), \(q_0\)만 보는 것이 아니라, PINN이 공간과 시간 전체에서 FDM 기준장을 얼마나 잘 따라가는지 확인한다.

\section{노트북 전체 흐름}

\begin{center}
\begin{tabular}{cll}
\toprule
셀 & 종류 & 역할 \\
\midrule
1 & Markdown & 노트북 목적과 비교 항목 설명 \\
2 & Code & import, 환경변수, 출력 폴더 설정 \\
3 & Markdown & 학습 설정 설명 \\
4 & Code & iteration, network, collocation point 수 설정 \\
5 & Markdown & PINN 학습 단계 설명 \\
6 & Code & DeepXDE PINN 학습 및 inverse parameter 출력 \\
7 & Markdown & FDM 기준장 생성 단계 설명 \\
8 & Code & explicit FDM으로 full-field reference 생성 \\
9 & Markdown & PINN 예측 단계 제목 \\
10 & Code & PINN을 FDM과 같은 격자에서 평가 \\
11 & Markdown & error metric 단계 제목 \\
12 & Code & MAE, RMSE, max error, relative L2 계산 \\
13 & Markdown & contour 비교 단계 제목 \\
14 & Code & FDM, PINN, absolute error contour plot 저장 \\
15 & Markdown & time slice 비교 단계 제목 \\
16 & Code & 특정 시간의 \(X-\theta\) profile 비교 \\
17 & Markdown & sensor history 비교 단계 제목 \\
18 & Code & sensor 위치별 시간 이력 비교 \\
19 & Markdown & 보고서용 요약 단계 제목 \\
20 & Code & 최종 수치 결과 출력 \\
\bottomrule
\end{tabular}
\end{center}

\section{셀 1: 제목과 목적 설명}

첫 번째 셀은 Markdown 셀이다.
노트북이 하는 일을 다음 네 가지로 정의한다.

\begin{itemize}[leftmargin=*]
    \item FDM 기준장 contour
    \item PINN 예측장 contour
    \item absolute error contour
    \item 시간 slice와 sensor 위치 time history 비교
\end{itemize}

여기서 ``full-field''라는 말은 sensor point 몇 개만 비교하지 않고, 모든 \(X\) 위치와 모든 시간 \(\tau\)에서 \(\theta\) 값을 비교한다는 뜻이다.

\section{셀 2: import와 실행 환경 설정}

두 번째 셀은 노트북 실행에 필요한 package와 project module을 불러온다.

\begin{lstlisting}[language=Python]
os.environ.setdefault("DDE_BACKEND", "pytorch")
os.environ.setdefault("PINN_FORCE_CPU", "1")
os.environ.setdefault("MPLCONFIGDIR", "/private/tmp/matplotlib-cache")
\end{lstlisting}

\begin{itemize}[leftmargin=*]
    \item \texttt{DDE\_BACKEND}: DeepXDE backend를 PyTorch로 고정한다.
    \item \texttt{PINN\_FORCE\_CPU}: M1 Mac에서 MPS/Metal crash를 피하기 위해 CPU 실행을 기본값으로 둔다.
    \item \texttt{MPLCONFIGDIR}: Matplotlib cache 위치를 writable temp directory로 지정한다.
\end{itemize}

그 다음 \texttt{main\_deepxde.py}를 \texttt{m}이라는 이름으로 import한다.
이 노트북은 물리 상수, FDM synthetic data 설정, DeepXDE model builder를 직접 다시 정의하지 않고 \texttt{main\_deepxde.py}의 값을 재사용한다.

\begin{lstlisting}[language=Python]
import main_deepxde as m
\end{lstlisting}

출력 그림은 다음 폴더에 저장된다.

\begin{lstlisting}[language=Python]
OUTPUT_DIR = Path("outputs/fdm_vs_pinn")
OUTPUT_DIR.mkdir(parents=True, exist_ok=True)
\end{lstlisting}

마지막 print들은 실행 환경과 true parameter를 확인하기 위한 것이다.
특히 \(\tau_{\mathrm{final}}\), \(\beta_{\mathrm{true}}\), \(S_{\mathrm{true}}\)가 기대한 값인지 확인한다.

\section{셀 3: 학습 설정 안내}

세 번째 셀은 Markdown 셀이다.
노트북 기본값이 빠른 smoke test용임을 설명한다.
보고서용으로는 다음 값을 올리는 것이 적절하다.

\begin{center}
\begin{tabular}{ll}
\toprule
설정값 & 보고서용 예시 \\
\midrule
\texttt{ITERATIONS} & \texttt{100000} \\
\texttt{NUM\_DOMAIN} & \texttt{2560} \\
\texttt{NUM\_BOUNDARY} & \texttt{256} \\
\texttt{NUM\_INITIAL} & \texttt{128} \\
\texttt{NUM\_TEST} & \texttt{1000} \\
\texttt{FDM\_N\_TIME} & \texttt{201} \\
\bottomrule
\end{tabular}
\end{center}

\section{셀 4: 학습 configuration 생성}

네 번째 셀은 실제 학습 설정값을 만든다.

\begin{lstlisting}[language=Python]
ITERATIONS = 1000
DISPLAY_EVERY = 100
SEED = 0

NUM_DOMAIN = 512
NUM_BOUNDARY = 64
NUM_INITIAL = 64
NUM_TEST = 200

FDM_NX = 101
FDM_N_TIME = 101
\end{lstlisting}

\texttt{ITERATIONS}는 Adam optimizer를 몇 step 돌릴지 정한다.
\texttt{DISPLAY\_EVERY}는 loss와 parameter를 몇 iteration마다 출력할지 정한다.
\texttt{FDM\_NX}와 \texttt{FDM\_N\_TIME}은 FDM 기준장을 저장할 공간/시간 해상도이다.

\begin{lstlisting}[language=Python]
args = Namespace(
    iterations=ITERATIONS,
    lr=1.0e-3,
    depth=3,
    width=32,
    num_domain=NUM_DOMAIN,
    num_boundary=NUM_BOUNDARY,
    num_initial=NUM_INITIAL,
    num_test=NUM_TEST,
    sensor_time_points=41,
    display_every=DISPLAY_EVERY,
    lbfgs=False,
)
\end{lstlisting}

\texttt{args}는 command line argument처럼 생긴 설정 묶음이다.
\texttt{main\_deepxde.py}의 \texttt{build\_deepxde\_model(args)}가 이 객체를 받아서 network, 학습점 개수, sensor data 개수 등을 설정한다.

\begin{lstlisting}[language=Python]
m.dde.config.set_random_seed(SEED)
\end{lstlisting}

random seed를 고정하면 같은 설정에서 결과 재현성이 좋아진다.

\section{셀 5: PINN 학습 단계 안내}

다섯 번째 셀은 Markdown 셀이다.
이 셀은 실제 학습 구현이 \texttt{main\_deepxde.py}의 \texttt{build\_deepxde\_model()}을 사용한다는 점을 설명한다.

즉, 노트북은 별도의 PINN model을 새로 만드는 것이 아니라 이미 정리된 DeepXDE model builder를 호출한다.

\section{셀 6: DeepXDE PINN 학습}

여섯 번째 셀은 PINN을 실제로 학습한다.

\begin{lstlisting}[language=Python]
model, raw_beta, raw_s = m.build_deepxde_model(args)
\end{lstlisting}

\texttt{build\_deepxde\_model(args)}는 세 가지를 반환한다.

\begin{center}
\begin{tabular}{ll}
\toprule
반환값 & 의미 \\
\midrule
\texttt{model} & DeepXDE model 객체 \\
\texttt{raw\_beta} & 학습되는 raw inverse parameter \\
\texttt{raw\_s} & 학습되는 raw inverse parameter \\
\bottomrule
\end{tabular}
\end{center}

\texttt{raw\_beta}와 \texttt{raw\_s}를 그대로 물리 파라미터로 쓰지 않는다.
실제 \(\beta\), \(S\)는 softplus를 통과시켜 양수로 만든다.

\[
\beta = \mathrm{softplus}(\texttt{raw\_beta}),
\qquad
S = \mathrm{softplus}(\texttt{raw\_s})
\]

다음 코드는 Adam optimizer로 model을 compile한다.

\begin{lstlisting}[language=Python]
model.compile(
    "adam",
    lr=args.lr,
    loss_weights=m.LOSS_WEIGHTS,
    external_trainable_variables=[raw_beta, raw_s],
)
\end{lstlisting}

\texttt{external\_trainable\_variables}에 \texttt{raw\_beta}, \texttt{raw\_s}를 넣어야 network weight뿐 아니라 inverse parameter도 같이 학습된다.

\begin{lstlisting}[language=Python]
losshistory, train_state = model.train(
    iterations=args.iterations,
    display_every=args.display_every,
    callbacks=[param_printer],
)
\end{lstlisting}

\texttt{param\_printer} callback은 학습 중 \(\beta\)와 \(S\)가 어떻게 변하는지 출력한다.
학습 후에는 다음 변환을 통해 실제 물리값으로 바꾼다.

\[
h = \beta \frac{kA}{PL^2}
\]

\[
q_0 = S \frac{k \Delta T_{\mathrm{ref}}}{I_{\mathrm{ref}}^2 L^2}
\]

\section{셀 7: FDM 기준장 생성 단계 안내}

일곱 번째 셀은 Markdown 셀이다.
PINN과 비교할 기준 해를 sensor 위치만이 아니라 전체 \(X-\tau\) 격자에서 만든다는 점을 설명한다.

이 단계에서 생성되는 기준장은 다음 배열이다.

\[
\theta_{\mathrm{FDM}}[n_{\tau}, n_X]
\]

\section{셀 8: explicit FDM full-field solver}

여덟 번째 셀은 이 노트북의 핵심 비교 기준을 만든다.
푸는 무차원 PDE는 다음과 같다.

\[
\frac{\partial \theta}{\partial \tau}
=
\frac{\partial^2 \theta}{\partial X^2}
- \beta_{\mathrm{true}} \theta
+ S_{\mathrm{true}} i(\tau)^2 g(X)
\]

\subsection*{격자 생성}

\begin{lstlisting}[language=Python]
x_values = np.linspace(0.0, 1.0, nx, dtype=np.float64)
tau_values = np.linspace(0.0, m.TAU_FINAL, n_time, dtype=np.float64)
\end{lstlisting}

\texttt{x\_values}는 \(0 \le X \le 1\)의 공간 격자이다.
\texttt{tau\_values}는 \(0 \le \tau \le \tau_{\mathrm{final}}\)의 저장용 시간 격자이다.

\subsection*{FDM 시간 간격}

\begin{lstlisting}[language=Python]
dx = x_values[1] - x_values[0]
dtau = 0.4 * dx**2
n_steps = math.ceil(m.TAU_FINAL / dtau)
\end{lstlisting}

explicit FDM에서는 시간 간격이 너무 크면 발산할 수 있다.
그래서 diffusion stability를 고려해 \(\Delta \tau = 0.4 \Delta X^2\)를 사용한다.

\subsection*{저장 배열}

\begin{lstlisting}[language=Python]
theta = np.zeros(nx, dtype=np.float64)
theta_xx = np.empty_like(theta)
theta_map = np.zeros((n_time, nx), dtype=np.float64)
\end{lstlisting}

\texttt{theta}는 현재 시간 step의 온도장이다.
\texttt{theta\_xx}는 공간 2차 미분 근사값이다.
\texttt{theta\_map}은 비교와 plot을 위해 저장되는 전체 \(X-\tau\) field이다.

\subsection*{열원과 true parameter}

\begin{lstlisting}[language=Python]
g_x = np.exp(-((x_values - m.X0) ** 2) / (2.0 * m.SIGMA_X**2))
source_on = float(m.S_TRUE) * g_x
beta_true = float(m.BETA_TRUE)
\end{lstlisting}

\(g(X)\)는 Gaussian heat source이다.
전류가 켜진 구간에서는 \(S_{\mathrm{true}}g(X)\)가 source term으로 들어가고, 꺼진 구간에서는 source term이 0이 된다.

\subsection*{단열 경계조건}

코드의 양 끝 경계는 다음과 같이 처리한다.

\begin{lstlisting}[language=Python]
theta_xx[0] = 2.0 * (theta[1] - theta[0]) / dx**2
theta_xx[-1] = 2.0 * (theta[-2] - theta[-1]) / dx**2
\end{lstlisting}

이는 \(X=0\), \(X=1\)에서

\[
\frac{\partial \theta}{\partial X}=0
\]

인 단열 경계조건을 ghost point 방식으로 반영한 형태이다.

\subsection*{시간 전진}

\begin{lstlisting}[language=Python]
theta = theta + step_dtau * (theta_xx - beta_true * theta + source)
\end{lstlisting}

이 식은 explicit Euler time stepping이다.
현재 온도장 \(\theta^n\)에서 다음 온도장 \(\theta^{n+1}\)을 직접 계산한다.

\subsection*{출력}

함수는 다음 세 개를 반환한다.

\begin{center}
\begin{tabular}{ll}
\toprule
반환값 & 의미 \\
\midrule
\texttt{x\_values} & 공간 격자 \\
\texttt{tau\_values} & 시간 격자 \\
\texttt{theta\_map} & FDM 기준 온도장 \\
\bottomrule
\end{tabular}
\end{center}

마지막으로 실제 시간 plot을 위해 무차원 시간 \(\tau\)를 초 단위 시간으로 다시 바꾼다.

\begin{lstlisting}[language=Python]
time_seconds = tau_values * m.L**2 / m.ALPHA
\end{lstlisting}

\section{셀 9: PINN 예측 단계 안내}

아홉 번째 셀은 Markdown 셀이다.
FDM과 같은 \(X-\tau\) grid에서 PINN 예측값을 계산한다는 단계 제목이다.

\section{셀 10: PINN을 같은 격자에서 평가}

열 번째 셀은 FDM이 저장된 격자와 동일한 위치에서 PINN을 평가한다.

\begin{lstlisting}[language=Python]
TAU, X = np.meshgrid(tau_values, x_values, indexing="ij")
xtau_eval = np.column_stack([X.reshape(-1), TAU.reshape(-1)]).astype("float32")
\end{lstlisting}

\texttt{meshgrid}는 \(X\)와 \(\tau\)의 2D grid를 만든다.
\texttt{model.predict()}는 입력을 \((N,2)\) 형태로 받기 때문에, 2D grid를 flatten해서 \texttt{xtau\_eval}로 바꾼다.

\begin{lstlisting}[language=Python]
theta_pinn = model.predict(xtau_eval).reshape(len(tau_values), len(x_values))
\end{lstlisting}

예측 후에는 다시 \((n_{\tau}, n_X)\) 배열로 reshape한다.
이렇게 해야 \texttt{theta\_fdm}과 같은 shape이 되어 직접 뺄 수 있다.

\section{셀 11: error metric 단계 안내}

열한 번째 셀은 Markdown 셀이다.
다음 셀에서 full-field error metric을 계산한다는 제목이다.

\section{셀 12: full-field error metric 계산}

열두 번째 셀은 PINN 예측장과 FDM 기준장의 차이를 계산한다.

\begin{lstlisting}[language=Python]
error = theta_pinn - theta_fdm
abs_error = np.abs(error)
\end{lstlisting}

계산하는 metric은 네 개다.

\begin{center}
\begin{tabular}{ll}
\toprule
Metric & 의미 \\
\midrule
MAE & 평균 절대 오차 \\
RMSE & root mean square error \\
Max error & 전체 격자에서 가장 큰 절대 오차 \\
Relative L2 & FDM field norm 대비 error field norm \\
\bottomrule
\end{tabular}
\end{center}

Relative L2는 다음과 같다.

\[
\mathrm{Relative\ L2}
=
\frac{\|\theta_{\mathrm{PINN}}-\theta_{\mathrm{FDM}}\|_2}
{\|\theta_{\mathrm{FDM}}\|_2}
\]

\section{셀 13: contour 비교 단계 안내}

열세 번째 셀은 Markdown 셀이다.
FDM, PINN, absolute error를 contour plot으로 비교한다는 제목이다.

\section{셀 14: full-field contour plot}

열네 번째 셀은 세 개의 contour plot을 한 figure에 그린다.

\begin{lstlisting}[language=Python]
fig, axes = plt.subplots(1, 3, figsize=(18, 4.5), constrained_layout=True)
\end{lstlisting}

세 subplot의 의미는 다음과 같다.

\begin{center}
\begin{tabular}{ll}
\toprule
Plot & 의미 \\
\midrule
FDM reference & FDM으로 계산한 기준 온도장 \\
PINN prediction & PINN이 예측한 온도장 \\
Absolute error & \(|\theta_{\mathrm{PINN}}-\theta_{\mathrm{FDM}}|\) \\
\bottomrule
\end{tabular}
\end{center}

그림은 다음 파일로 저장된다.

\begin{lstlisting}[language=Python]
outputs/fdm_vs_pinn/fdm_pinn_full_field_contours.png
\end{lstlisting}

보고서에서는 이 그림이 전체 field 비교의 핵심 figure가 된다.

\section{셀 15: time slice 비교 단계 안내}

열다섯 번째 셀은 Markdown 셀이다.
특정 시간에서 \(X\) 방향 온도 profile을 비교한다는 제목이다.

\section{셀 16: selected time slice 비교}

열여섯 번째 셀은 네 개 시간에서 FDM과 PINN의 \(X-\theta\) profile을 비교한다.

\begin{lstlisting}[language=Python]
slice_times = [300.0, 600.0, 900.0, 1200.0]
\end{lstlisting}

각 시간은 current profile의 전류 on/off 전환과 관련 있다.

\begin{center}
\begin{tabular}{ll}
\toprule
시간 & 의미 \\
\midrule
300 s & 첫 번째 heating 구간 끝 \\
600 s & 첫 번째 cooling 구간 끝 \\
900 s & 두 번째 heating 구간 끝 \\
1200 s & 전체 해석 시간 끝 \\
\bottomrule
\end{tabular}
\end{center}

\begin{lstlisting}[language=Python]
idx = int(np.argmin(np.abs(time_seconds - target_time)))
\end{lstlisting}

위 코드는 원하는 실제 시간에 가장 가까운 FDM 저장 index를 찾는다.
그 다음 같은 index에서 FDM과 PINN profile을 함께 plot한다.

그림은 다음 파일로 저장된다.

\begin{lstlisting}[language=Python]
outputs/fdm_vs_pinn/fdm_pinn_time_slices.png
\end{lstlisting}

\section{셀 17: sensor history 비교 단계 안내}

열일곱 번째 셀은 Markdown 셀이다.
sensor 위치에서 시간에 따른 온도 이력을 비교한다는 제목이다.

\section{셀 18: sensor 위치별 time history 비교}

열여덟 번째 셀은 sensor 위치 \(X=0.25, 0.55, 0.85\)에서 FDM과 PINN의 시간 이력을 비교한다.

\begin{lstlisting}[language=Python]
for ax, sensor_x in zip(axes, m.SENSOR_X):
    idx = int(round(sensor_x * (len(x_values) - 1)))
\end{lstlisting}

\texttt{sensor\_x}는 무차원 위치이므로, \texttt{len(x\_values)-1}을 곱해 가장 가까운 grid index로 변환한다.

\begin{lstlisting}[language=Python]
ax.plot(time_seconds, theta_fdm[:, idx], label="FDM", linewidth=2)
ax.plot(time_seconds, theta_pinn[:, idx], "--", label="PINN", linewidth=2)
\end{lstlisting}

이 plot은 기존 inverse PINN 학습에서 사용한 sparse sensor 위치에서 두 해가 얼마나 잘 맞는지 보여준다.
그림은 다음 파일로 저장된다.

\begin{lstlisting}[language=Python]
outputs/fdm_vs_pinn/fdm_pinn_sensor_histories.png
\end{lstlisting}

\section{셀 19: 보고서용 요약 단계 안내}

열아홉 번째 셀은 Markdown 셀이다.
마지막 셀에서 보고서에 옮기기 쉬운 형태로 주요 수치를 출력한다는 제목이다.

\section{셀 20: 최종 summary 출력}

마지막 셀은 학습 설정, inverse parameter, error metric을 한 번에 출력한다.

\begin{lstlisting}[language=Python]
print(f"iterations      : {ITERATIONS}")
print(f"network         : {args.depth} x {args.width}")
print(f"loss weights    : {m.LOSS_WEIGHTS}")
\end{lstlisting}

학습 설정을 기록하는 이유는 결과가 iteration 수, network 크기, loss weight에 따라 달라지기 때문이다.

\begin{lstlisting}[language=Python]
print(f"beta learned    : {beta_learned:.6f} / true {m.BETA_TRUE:.6f}")
print(f"S learned       : {s_learned:.6f} / true {m.S_TRUE:.6f}")
print(f"h learned       : {h_learned:.6f} / true {m.H_TRUE:.6f}")
print(f"q0 learned      : {q0_learned:.6e} / true {m.Q0_TRUE:.6e}")
\end{lstlisting}

이 부분은 inverse problem이 true parameter를 얼마나 잘 복원했는지 보여준다.

\begin{lstlisting}[language=Python]
print(f"MAE             : {mae:.6e}")
print(f"RMSE            : {rmse:.6e}")
print(f"Max error       : {max_error:.6e}")
print(f"Relative L2     : {rel_l2:.6e}")
\end{lstlisting}

이 부분은 full-field prediction accuracy를 수치로 정리한다.

\section{보고서에 쓸 때 핵심 포인트}

이 노트북 결과를 보고서에 넣을 때는 다음 세 가지를 함께 제시하는 것이 좋다.

\begin{enumerate}[leftmargin=*]
    \item learned inverse parameter와 true parameter 비교
    \item FDM/PINN/error contour figure
    \item MAE, RMSE, relative L2 같은 full-field error metric
\end{enumerate}

단순히 \(h\), \(q_0\)가 true value와 비슷하다고 끝내면 부족하다.
PINN이 PDE solution field 자체도 잘 재현하는지 보여줘야 하므로 full-field comparison이 필요하다.

\end{document}
